% =====================================================================
%  SPRACH-SCHALTER der Diplomarbeit  (bilingual single source)
%  Werte:  de  (Deutsch, primär)   |   en  (Englisch, sekundär)
%
%  Diese eine Zeile umschalten, um die Sprache zu wechseln (z. B. auf
%  Overleaf, wo kein Kommandozeilen-Override möglich ist).
%  Lokal kann der Wert per Build-Skript überschrieben werden:
%      pwsh -File build.ps1 -Lang de
%  (siehe sessions/2026-06-01-bilingualer-compile-schalter.md)
% =====================================================================
\providecommand{\thesislang}{de}

% ---------------------------------------------------------------------
%  [EN-1, 22.09.2026] Sprachabhaengige Zugriffsnotiz der Literatur-Eintraege
%  (literatur.bib, Felder howpublished/note): DE 'abgerufen am', EN 'accessed'.
%  \ifthesisDE wird in diplomarbeit.tex aus \thesislang abgeleitet; das Makro
%  wird erst im Literaturverzeichnis expandiert, also nach dieser Ableitung.
%  [FIX-EN r1, 22.09.2026] Erweiterung (LENS-EN-A S-1/S-2, LENS-EN-B S17/S21):
%    \biblang{DE-Text}{EN-Text}  sprachabhaengiger Feld-/Titeltext eines Eintrags
%    \bibdate{TT.MM.JJJJ}         DE unveraendert TT.MM.JJJJ, EN ISO JJJJ-MM-TT
%    \bibaccessed{TT.MM.JJJJ}     DE 'abgerufen am TT.MM.JJJJ', EN 'accessed JJJJ-MM-TT'
%  Gebrauch in literatur.bib:  \bibaccessed{19.09.2026}
% ---------------------------------------------------------------------
\providecommand{\biblang}[2]{\ifthesisDE{#1}\else{#2}\fi}
\providecommand{\bibdate}[1]{\bibdateparse#1\bibdatestop}
\def\bibdateparse#1.#2.#3\bibdatestop{\ifthesisDE{#1.#2.#3}\else{\mbox{#3-#2-#1}}\fi}
\providecommand{\bibaccessed}[1]{\biblang{abgerufen am}{accessed} \bibdate{#1}}
